\begin{ExerciseBlock}
\ExerciseHeading{06}{Mehrdimensionaler inelastischer Stoß bei Autounfall}{two_particle_systems_two_dimensional_car_collision}
Ein Fahrzeug $m_1 = 1,0$\,t fährt mit $v_1 = 100\,$km/h auf einer Landstraße entlang. Senkrecht dazu fährt ein Auto ($m_2 = 1,2\,$t) mit $v_2 = 40\,$km/h ein. Es kommt zur Kollision, bei der sich beide Autos vollkommen ineinander verkeilen und zusammen mit der Geschwindigkeit $\vec{v}$ und im Winkel $alpah$ zur Fahrtrichtung von $m_1$ von der Fahrbahn fliegen (siehe Abbildung).

\begin{center}
	\CachedTikzFigure{content/resources/chapters/two_particle_systems/tikz/two_dimensional_car_collision.tex}
\end{center}

\begin{enumerate}
	\item Verwenden Sie die Impulserhaltung, um die Geschwindigkeit $v$ sowie die Winkel $\alpha$ zu bestimmen. Bedenken Sie, dass der Impuls eine vektorielle Größe ist.

	\item Berechnen Sie die Energie $Q$, die in Verformung und Wärme umgewandelt wurde.

	\item Wie groß ist der Anteil von $Q$ an der ursprünglichen Bewegungsenergie beider Fahrzeuge.

\end{enumerate}
\end{ExerciseBlock}
\UseExerciseSolution{two_particle_systems_two_dimensional_car_collision}
